% !TeX root = main.tex
% Engine: XeLaTeX, selected by scripts/build.py.

\documentclass[11pt, aspectratio=169]{beamer}
\input{build/toolkit-path.tex}
\input{\KGUPCToolkitRoot/latex/layouts/beamer.tex}
\hypersetup{filecolor=solutio-purple}
\usepackage{catchfile}
\usepackage{xr}
\usetikzlibrary{arrows.meta,matrix,decorations.pathreplacing}

% 어드밴스드를 먼저 빌드하면 PDF의 실제 페이지 번호를 가져올 수 있다.
\externaldocument[advanced-]{../../2026-advanced/solutions/main}
\newcommand{\AdvancedEditorialURL}{https://kgu-solutio.github.io/kgupc-archive/2026-advanced/solutions/main.pdf}
\newcommand{\advancededitoriallink}[2]{%
    \ifnum\getpagerefnumber{advanced-#1}=0\relax
        \PackageError{editorial-links}{Advanced editorial page is missing: #1}
            {Build 2026-advanced/solutions/main.tex before the beginner editorial.}%
    \fi
    \href{\AdvancedEditorialURL\#page=\getpagerefnumber{advanced-#1}}{#2}%
}

% 문제 목록의 배치를 유지하고, 아래쪽 빈 공간에 후기 링크를 표시한다.
\newcommand{\afterwordlink}{%
    \begin{tikzpicture}[remember picture,overlay]
        \node[anchor=south,inner sep=0pt,font=\footnotesize]
            at ([yshift=1cm]current page.south)
            {\hyperlink{afterword}{출제진 후기 보기}};
    \end{tikzpicture}%
}

\title[KGUPC 2026 (Beginner) 에디토리얼]
      {KGUPC 2026 (Beginner) \newline 에디토리얼}
\subtitle{Official Solutions}
\author{Solutio in Kyonggi Univ. KGUPC 2026 출제진}
% 기존 2025 템플릿의 이미지. 2026 후원사 이미지로 교체한다.
\renewcommand{\sponsorimage}{../sponsors.png}
\date{2026년 10월 30일}

% 문제 제목은 원본 name.tex에서 읽는다. 제목을 중복으로 수정하지 않는다.
\CatchFileDef{\ProblemTitleA}{../problems/A/statement-sections/korean/name.tex}{\endlinechar=-1}
\CatchFileDef{\ProblemTitleB}{../problems/B/statement-sections/korean/name.tex}{\endlinechar=-1}
\CatchFileDef{\ProblemTitleC}{../problems/C/statement-sections/korean/name.tex}{\endlinechar=-1}
\CatchFileDef{\ProblemTitleD}{../problems/D/statement-sections/korean/name.tex}{\endlinechar=-1}
\CatchFileDef{\ProblemTitleE}{../problems/E/statement-sections/korean/name.tex}{\endlinechar=-1}
\CatchFileDef{\ProblemTitleF}{../problems/F/statement-sections/korean/name.tex}{\endlinechar=-1}
\CatchFileDef{\ProblemTitleG}{../problems/G/statement-sections/korean/name.tex}{\endlinechar=-1}
\CatchFileDef{\ProblemTitleH}{../problems/H/statement-sections/korean/name.tex}{\endlinechar=-1}
\CatchFileDef{\ProblemTitleI}{../problems/I/statement-sections/korean/name.tex}{\endlinechar=-1}
\CatchFileDef{\ProblemTitleJ}{../problems/J/statement-sections/korean/name.tex}{\endlinechar=-1}
\CatchFileDef{\ProblemTitleK}{../problems/K/statement-sections/korean/name.tex}{\endlinechar=-1}
\CatchFileDef{\ProblemTitleL}{../problems/L/statement-sections/korean/name.tex}{\endlinechar=-1}

\begin{document}
    \setcounter{framenumber}{-1}
    \frame{\titlepage}

    \begin{frame}{문제 목록}
        \begin{center}
            \setlength{\tabcolsep}{4.5pt}
            \begin{tabular}{cl|c|c}
                \hline
                문제 & & 예상 난이도 & 출제자 \\
                \hline
                \problemlink{A}{\textbf{A}} & \problemlink{A}{\ProblemTitleA} & \tiertext{Bronze4} & areian13 \\
                \problemlink{B}{\textbf{B}} & \problemlink{B}{\ProblemTitleB} & \tiertext{Bronze3} & areian13 \\
                \problemlink{C}{\textbf{C}} & \problemlink{C}{\ProblemTitleC} & \tiertext{Bronze3} & mingmingmon \\
                \problemlink{D}{\textbf{D}} & \problemlink{D}{\ProblemTitleD} & \tiertext{Bronze2} & mingmingmon \\
                \problemlink{E}{\textbf{E}} & \problemlink{E}{\ProblemTitleE} & \tiertext{Bronze2} & areian13 \\
                \problemlink{F}{\textbf{F}} & \problemlink{F}{\ProblemTitleF} & \tiertext{Bronze2} & areian13 \\
                \problemlink{G}{\textbf{G}} & \problemlink{G}{\ProblemTitleG} & \tiertext{Bronze2} & areian13 \\
                \problemlink{H}{\textbf{H}} & \problemlink{H}{\ProblemTitleH} & \tiertext{Bronze1} & mingmingmon \\
                \problemlink{I}{\textbf{I}} & \problemlink{I}{\ProblemTitleI} & \tiertext{Bronze1} & mingmingmon \\
                \problemlink{J}{\textbf{J}} & \problemlink{J}{\ProblemTitleJ} & \tiertext{Silver3} & areian13 \\
                \hline
            \end{tabular}
        \end{center}
    \end{frame}

    \begin{frame}{문제 목록}
        \begin{center}
            \setlength{\tabcolsep}{4.5pt}
            \begin{tabular}{cl|c|c}
                \hline
                문제 & & 예상 난이도 & 출제자 \\
                \hline
                \problemlink{K}{\textbf{K}} & \problemlink{K}{\ProblemTitleK} & \tiertext{Silver4} & ehdnjs14 \\
                \problemlink{L}{\textbf{L}} & \problemlink{L}{\ProblemTitleL} & \tiertext{Silver3} & areian13 \\
                \hline
            \end{tabular}%
            \afterwordlink
        \end{center}
    \end{frame}

    % 풀이 순서는 문제셋과 동일하다.
    \import{A}{from-npown-to-1-small.tex}
    \import{B}{my-name-is-k1ry0n6.tex}
    \import{C}{kiryong-is-hangry.tex}
    \import{D}{kiryong-budget-book.tex}
    \import{E}{every-monent-saving-life.tex}
    \import{F}{do-you-need-multitap.tex}
    \import{G}{one-more-picture-plz.tex}
    \import{H}{kiryong-loves-movie.tex}
    \import{I}{to-do-list.tex}
    \import{J}{what-is-my-pair.tex}
    \import{K}{ice-water-purifier-small.tex}
    \import{L}{cut-dujjonku.tex}
    % !TeX root = main.tex

% 이번 대회는 두 부문의 후기를 어드밴스드 에디토리얼에 함께 작성한다.
\section{출제진 후기}

\begin{frame}[label=afterword]{출제진 후기}
    \begin{center}
        출제진 후기는 어드밴스드 에디토리얼에 함께 모았습니다.

        \vspace{5mm}
        \advancededitoriallink{afterword}{어드밴스드 출제진 후기 보기}
    \end{center}
\end{frame}

\end{document}
